\subsection{Round seven: a placement corpus and a hindsight clustering}
\label{sec:r7-method}
Round seven asks two questions of data instead of testing an idea (N Q1.7: how far is the VPR packer from a best-known clustering made with
hindsight from good placements, and how much would LUT duplication add; N Q2.4: which HDL relations predict that atoms sit together in good
placements). Its protocol, committed before any run, is \code{docs/plan/round7.md} (commit \code{93c983a}, which froze the text only; the run code was
committed later, as \code{94827f5}); the one quantity defined in
advance is the \emph{hindsight gap} per design, and every other analysis is exploratory.

\textbf{Corpus.} 15 VTR designs of rounds one to four and 12 round-five netlist-level constructions
(mesh and tree, $B\in\{16,64,256\}$, whose block clustering is known by construction; it is a known clustering, not an optimum), each with VPR seeds 1 to 10 (stock and affinity-bias runs) or 1 to 5 (\code{.net} arms), the architecture of rounds one
to six, VPR's default analytical-placement flow for the stock runs (VTR~\cite{rose2012vtr} and VTR~9~\cite{elgammal2025vtr9}, commit \code{3c9a4d23b2}, version 9.0.0-dev; the flow already clusters with APPack from atom positions; there is no separate analytical hindsight
start), and routing at the minimum channel width. All designs of the main set are included. Each run records, per atom, the construct chain from the provenance names, the word and bit
index, the CLB, its coordinates, the pack-time (stock arm only) and post-route criticality and slack; per pair, the N17 class, the HDL relations and the distance in
each run; per net, fanout and CLBs spanned. Post-route criticality is path based, from the 200 worst reported paths.

\textbf{Hindsight clustering (the \code{.net} arm).} The stock clustering is read from the stock \code{.net} of each design (the same in every seed). Its unit is the logic element
(fle: a LUT with its flip-flop, or a pair of small LUTs); CLBs with carry chains stay as they are. Take the best run of the previous arm
(seeds 1, 2, 3, 4, 5; lowest wirelength~$\times$~critical-path delay), keep its CLB positions fixed, and move logic elements between CLBs within Manhattan distance two
(or swap them) by greedy search, minimising the sum over nets of $(1+4\,\mathrm{pcrit})\cdot$HPWL under the exact CLB legality of the architecture
(at most 10 logic elements, 40 distinct input nets and one clock per CLB; the crossbar is complete). The clustering is written as a \code{.net} file and VPR (VTR~\cite{rose2012vtr}, VTR~9~\cite{elgammal2025vtr9}, doi 10.1145/3734798, commit \code{3c9a4d23b2}) is run with packing off
(\code{--net\_file}, annealing placement from a random start, the router of every other arm), so the clustering VPR places is the one written (fidelity 1.0 by construction, asserted for every run
from the placed CLBs). \code{net0} is the stock clustering written by the same code and placed the same way, so \code{net0} and the hindsight arm differ only in the clustering. Repeat from
the best run of each round until the median objective stops improving by 1\% or after five rounds. The rule is fixed in advance: the arm is the round with the lowest median objective over the five seeds, and the \emph{gap} of a metric is
\code{net0} over that arm, seed for seed, with a bootstrap interval over those seeds. The gap on these seeds is \emph{in sample} (the arm is chosen on the seeds it is evaluated on; winner's-curse risk) and is labelled \emph{selection seeds}. The headline is the \emph{held-out} gap: after the arm is fixed on seeds 1, 2, 3, 4, 5, \code{net0} and the arm are run on seeds 6, 7, 8, 9, 10 (protocol amendment of 2026-10-27, written before the first held-out run), for every real design where an arm other than \code{net0} was chosen and, on the constructions, also for the known clustering. A design has a held-out gain when the 95\% bootstrap interval of its objective ratio (wirelength~$\times$~critical-path delay, \code{net0} over the arm) lies above 1. Peter cut the seeds of the new runs from ten to five (decision 2026-10-26, before the \code{.net} runs); the affinity-bias data have ten seeds and get an out-of-sample check (arm chosen on seeds 1 to 5, evaluated on seeds 6 to 10). Each \code{.net} file is also checked for equivalence with the BLIF and with the stock \code{.net} (every atom sees the same nets on its input pins; every net read outside its CLB has an output pin; \code{scripts/r7netcheck.py}). Where the chosen arm is \code{net0} the procedure found no gain; no gap is printed for the design, and the pooled medians count it as gap 1 (recovery 0) beside the medians over the designs with an arm. Stock packing is the
same in every seed, so the CLB-count gap has no clustering variance: its interval covers placement noise only.

\textbf{Affinity-bias arm.} The earlier arm of this round hands the clustering to the Norns packer as an affinity file at weight 100 on the analytical-placement flow. VPR's packer
rebuilds only part of the requested groups, so its gap is the \emph{gap reachable via affinity bias}, not a lower bound on the gap between the packer and a best-known clustering; its fidelity is reported.
\emph{Duplication} candidates are single LUTs whose output net has sinks in two or more logic CLBs; nothing is built.
